\documentclass[11pt]{article}

\usepackage[margin=1in]{geometry}
\usepackage{amsmath,amssymb,amsthm}
\usepackage{mathtools}
\usepackage{booktabs}
\usepackage{enumitem}
\usepackage[colorlinks=true,linkcolor=blue,citecolor=blue,urlcolor=blue]{hyperref}
\usepackage[numbers,sort&compress]{natbib}

\newtheorem{theorem}{Theorem}
\newtheorem{proposition}[theorem]{Proposition}
\newtheorem{lemma}[theorem]{Lemma}
\newtheorem{corollary}[theorem]{Corollary}
\newtheorem{definition}[theorem]{Definition}
\newtheorem{remark}[theorem]{Remark}
\newtheorem{obligation}{Proof Obligation}

\newcommand{\R}{\mathbb{R}}
\newcommand{\xv}{\mathbf{x}}
\newcommand{\pv}{\mathbf{p}}
\newcommand{\uv}{\mathbf{u}}
\newcommand{\Mopt}{\mathcal{M}_{\omega_0}}

\title{Chart-Aware Contiguous Memory for Discretized Optical Phase Space:\\
Geometry, Addressing, and the Limits of a Four-Day Prototype}

\author{Sethurathienam Iyer\thanks{ORCID: 0009-0008-5446-2856.
Independent work; associated Zenodo records are cited in the bibliography.}\\
BITS Pilani, Goa Campus, Sancoale, India}

\date{October 22, 2026}

\begin{document}
\maketitle

\begin{abstract}
We develop a chart-aware, contiguous-memory architecture for discretized optical
phase space, separating ray dynamics, measurement observables, geometric
discretization, and address correctness. A progressively ellipsoidal optical domain
motivates the use of Hamiltonian characteristics and local phase-space charts rather
than conventional Cartesian indexing. We derive the geometric conditions under which
the domain is well posed, show that spatially indexed memory loses ray-branch identity
at caustics, and formulate the address construction as a multi-chart encoder with
explicit range and injectivity invariants, released as a verification-ready Ada/SPARK
specification with a GNATprove transcript (47/47 proof obligations discharged, 0
unproved, for this encoder instance and its contract structure). Under a fixed hardware
budget we formalize cell allocation as a task-aware quantization problem via the law of
total variance, with calibration-effort costs estimated from finite samples. On a
released calibrated synthetic model, a chart-legal dynamic-programming partition
improves held-out timing distortion by $\mathbf{1.75\times}$ over uniform angular binning
($4.230\times10^{-4}$ vs $2.423\times10^{-4}$; ratio $1.746\pm0.005$ at 95\% CI, 30
sampling seeds). We report controls that bound what this number means: the gain falls
from $2.41\times$ to $1.51\times$ as the budget grows from $K=16$ to $K=256$; a single
smooth delay law with no two-regime structure already gives $1.50\times$; and because
the delay law is known in closed form, analytic quadrature matches the sample-estimated
partition to $0.07\%$. We further price the chart-legal constraint by computing what it
forbids: optimizing the scalar observable alone reaches $10.6\times$ lower distortion,
but its cells fragment into $91$ angular runs against $32$, so it is not
representable at equal address cost. Chart legality is therefore an
address-structure cost, not a moral one, and topology, identity, address count,
calibration, and distortion are coupled resources. Synthetic experiments demonstrate
that a spatial-only representation can manufacture apparent timing fluctuations by
merging distinct ray branches at a caustic, and that finite-resolution phase-space
representation converts between-branch variance into within-cell variance by an exactly
quantifiable identity.
The generator, split rule, partition code, and chart table are released so that these
figures are reproducible. We state the limits of what can be established in four days:
machine-checked verification of the complete encoder beyond the released
specification, broadband propagation accuracy, geometric approximation quality, and
metrological benefit remain explicit validation obligations.
\end{abstract}

\tableofcontents

\section{Introduction}\label{sec:intro}

Consider an intern tasked with designing a memory structure for measurements that live
in three spatial dimensions and two rotational dimensions, occupying a volume whose
geometry becomes progressively more elliptical toward the boundary, with contiguous
storage, in Ada, for detecting femtosecond-scale fluctuations in an atomic clock---on a
four-day deadline. Every noun in this assignment conceals another discipline: memory
conceals topology, ellipsoid conceals differential geometry, optics conceals
Hamiltonian mechanics, femtosecond conceals precision metrology, and Ada conceals
formal verification.

The purpose of this paper is not to discover a new clock. It is to explain why a wrong
memory layout can invent one. The surprising central idea is that the correct memory
addresses may need to follow the geometry of light propagation, rather than the
coordinates occupied by matter. The central claim is not that a five-dimensional array
is better than a three-dimensional array. It is that the optimal representation is not
the one that stores the fewest coordinates or minimizes a single timing error; it is
the least expensive representation that preserves every distinction the instrument may
need to make.

\subsection{The design problem, formally}\label{sec:problem}

The paper's main technical contribution can be stated as a single constrained
optimization problem before any of its components are derived. Let $Q$ be an optical
state drawn from the instrument's operating distribution $\mu$ on the fixed-frequency
manifold $\Mopt$; let $T_g(Q)$ be its group-delay correction; and let
$B_{\mathcal{P}}(Q)$ denote the cell of a partition $\mathcal{P}=\{C_1,\ldots,C_K\}$
containing $Q$. The design problem is
\begin{equation}
\boxed{
\begin{aligned}
\min_{\mathcal{P},\{m_k\}} \quad
& \mathbb{E}\!\left[\operatorname{Var}\!\left(T_g(Q)\mid B_{\mathcal{P}}(Q)\right)\right]
  + \sum_{k=1}^{K} \frac{p_k v_k}{m_k} \\[2pt]
\text{subject to}\quad
& |\mathcal{P}| \le K_{\max}, \\
& \textstyle\sum_{k=1}^{K} m_k \le M, \\
& \text{each cell has one canonical chart owner and one address,} \\
& \text{addresses are contiguous and collision-free (Theorem~\ref{thm:injective}),} \\
& \text{required ray-branch distinctions are preserved (Theorem~\ref{thm:branchcost}),} \\
& \text{the chart table is frozen per acquisition batch,}
\end{aligned}}
\end{equation}
where $p_k=\mathbb{P}(Q\in C_k)$, $v_k=\operatorname{Var}(T_g\mid C_k)$, $K_{\max}$ is
the address budget, and $M$ is the calibration budget. The objective's first term is
the geometric distortion resolved by the partition (exactly decomposed by the law of
total variance, \S\ref{sec:rate-distortion}); the second is the statistical cost of
estimating each cell's correction from $m_k$ calibration samples. The constraints
encode the architectural requirements that are not tradeable: representation
correctness (contiguity, uniqueness, frozen versioning) and preservation of the
optical distinctions the measurement needs. Everything that follows---the manifold
geometry, the encoder, the rate-distortion analysis, and the experiments---is in
service of this problem. Making it explicit up front also makes the paper's central
claim precise: the optimum is not the partition that minimizes timing distortion
alone; it is the cheapest partition in the feasible set.

\subsection{Evidence taxonomy}\label{sec:taxonomy}

Throughout, we keep three kinds of evidence distinct, because confusing them is how
prototypes become folklore.
\begin{enumerate}[label=(\Alph*)]
\item \textbf{Mathematical identities.} Proved: Snell's law from Fermat's principle,
the law of total variance, the range bound of the encoder.
\item \textbf{Synthetic numerical results.} Demonstrated under stated models with
stated sample sizes. These are not physical atomic-clock measurements; they establish
failure modes in controlled settings.
\item \textbf{Formally verified implementation claims.} These require a proof
transcript, e.g.\ a GNATprove transcript. \emph{A transcript is presented in
Appendix~\ref{app:ada} for the released encoder, its chart table, and its contract
structure only.} Exhaustive enumeration of a concrete chart table is evidence for that
table only; it is not a general proof, and the general statement remains
Theorem~\ref{thm:injective}.
\end{enumerate}
When we say ``proved'' we mean (A); when we say ``the simulation shows'' we mean (B);
when we say ``we have not proved,'' category (C) is missing. This distinction is the
paper's immune system.

\subsection{Related work}\label{sec:related}

The architectural roots of this problem are cross-disciplinary.
\emph{Optics and mechanics:} Fermat's principle and the eikonal equation are treated
authoritatively by Born and Wolf~\cite{bornwolf1999}; Hamiltonian characteristics,
phase-space flows, symplectic geometry, and Liouville's theorem by
Arnold~\cite{arnold1989}; caustics, singular ray projections, and Airy/Pearcey
catastrophe optics by Berry and Upstill~\cite{berry1980,pearcey1946,airy1841}; and the
contact-geometric structure of the fixed-frequency hypersurface by
Geiges~\cite{geiges2008}.
\emph{Metrology:} the statistics of atomic frequency standards under noise perturbations
whose character depends on sampling time are due to Allan~\cite{allan1966}; this grounds
the measurement model of \S\ref{sec:measurement} and the spectral diagnosis proposed in
\S\ref{sec:extensions}.
\emph{Information theory:} Shannon's coding theorems with a fidelity criterion provide the
rate-distortion foundation~\cite{shannon1959}; Lloyd~\cite{lloyd1982} and Max~\cite{max1960}
establish optimal scalar quantization; Gray and Neuhoff~\cite{gray1998} connect quantization,
distortion, resolution, and coding rate; and Misra, Goyal, and Varshney~\cite{misra2011} study
quantizers designed to minimize the error of a \emph{computed function}---the closest
established literature to our task-aware timing objective, with branch constraints and
calibration costs specifying the optical specialization. Stone~\cite{stone1982} supplies the
nonparametric estimation rates underlying the $M^{1/3}$ scaling law of
Theorem~\ref{thm:threshold}.
\emph{Optimization and verification:} the inverse-Hessian sensitivity of
Theorem~\ref{thm:dispersion} is an application of perturbation analysis of parameterized
optimization~\cite{bonnans2000}; and SPARK's contract-based verification methodology is
codified by McCormick and Chapin~\cite{mccormick2015}.
The paper also builds on the author's prior computational frameworks---physics-informed
combinatorial solvers using spectral, topological, and arithmetic methods
\cite{nitrosat,multiplicative,shunyabar,banzhaf,spectralmult,quantumvacuum}, sharing the
theme of phase-transition detection and branch-aware optimization---and earlier applied work
on content-based retrieval~\cite{imagecaptioning}; these are self-citations of record used
for context only, not as evidence of type (A), (B), or (C).

\subsection{Contributions}

\begin{enumerate}[label=(\roman*)]
\item A geometric model of progressively ellipsoidal domains with explicit nesting
conditions (\S\ref{sec:domain}), and the observation that compactness is a storage
property while locality is a geometric property, which are not interchangeable
(\S\ref{sec:arrays}).
\item Derivation of the optical state space from Fermat's principle, the eikonal
equation, and Hamiltonian characteristics; identification of caustics as the failure of
\emph{spatial} invertibility, not of ray dynamics (\S\ref{sec:fermat}--\S\ref{sec:phase-space}).
\item A chart-aware contiguous encoder with explicit range and injectivity obligations,
including the distinction between the two (\S\ref{sec:addressing}, \S\ref{sec:ada}),
released as a verification-ready specification with a GNATprove transcript
(Appendix~\ref{app:ada}).
\item Phase-versus-group-delay separation: the state manifold is generated by the phase
dispersion relation, while group delay is a different observable computed along its
rays (\S\ref{sec:femto}).
\item A task-aware rate-distortion formulation of cell allocation under memory and
calibration budgets (\S\ref{sec:rate-distortion}), and five synthetic experiments
demonstrating concrete failure modes, including a false clock manufactured by spatial
projection (\S\ref{sec:experiments}).
\item An explicit statement of what remains unproved: the SPARK transcript, the
geometric preprocessing quality, and the metrological benefit
(\S\ref{sec:limits}).
\end{enumerate}

\section{Why an ordinary array is the wrong starting point}\label{sec:arrays}

A conventional multidimensional array is a convenient interpretation of linear memory.
Let
\begin{equation}
(i,j,k), \qquad 0 \le i < N_x, \quad 0 \le j < N_y, \quad 0 \le k < N_z.
\end{equation}
A row-major representation uses
\begin{equation}
A(i,j,k) = (iN_y + j)N_z + k.
\end{equation}
This is simple, collision-free, and invertible. Adding two angular coordinates
discretized into $N_\theta$ and $N_\phi$ bins yields
\begin{equation}
A(i,j,k,p,q) = ((((iN_y+j)N_z+k)N_\theta+p)N_\phi+q.
\end{equation}
Five dimensions have become one integer. Problem solved? Not remotely.

This formula assumes that the logical domain is a Cartesian product
$X \times Y \times Z \times \Theta \times \Phi$. But an ellipsoidal volume is not a
rectangular box, and a progressively deforming ellipsoidal volume is not naturally
described by uniformly sampled Cartesian slices. Allocating the bounding box wastes
memory outside the valid domain; enumerating valid lattice points with consecutive
ranks saves memory but does not guarantee that neighboring optical states are stored
near one another.

A \emph{compact address map} answers: can I store every admissible state without
wasting addresses? An optical memory architecture must answer a harder question: can I
store states so that the geometry relevant to the measurement survives discretization?
Compactness is a storage property; locality is a geometric property. They are not
interchangeable.

Consider two measurements that are close in Euclidean space but arrive along different
ray branches with very different propagation delays. A compact spatial rank puts them
together; an optical memory that cares about timing should not necessarily. Conversely,
two measurements far apart in space but sharing the same group delay may be natural
comparison partners for timing analysis; a spatial layout separates them. The intern's
first mistake would be to make the array small rather than faithful.

\section{A progressively elliptical domain}\label{sec:domain}

Introduce a normalized radial parameter $0 \le \rho \le 1$ and a unit direction
$\uv \in S^2$. Define
\begin{equation}
\xv(\rho,\uv) = \rho
\begin{pmatrix} a(\rho) & 0 & 0 \\ 0 & b(\rho) & 0 \\ 0 & 0 & c(\rho) \end{pmatrix}
\uv,
\qquad
a(\rho) = 1 + \alpha\rho^2, \quad b(\rho) = 1, \quad c(\rho) = 1 - \beta\rho^2,
\end{equation}
with restrictions
\begin{equation}
\boxed{\alpha > 0, \qquad 0 < \beta < \tfrac13.}
\end{equation}
These are not cosmetic. The semiaxes of the shell at radial parameter $\rho$ are
\begin{equation}
s_x(\rho) = \rho + \alpha\rho^3, \qquad s_y(\rho) = \rho, \qquad s_z(\rho) = \rho - \beta\rho^3.
\end{equation}
For the ellipsoids to remain nested, their semiaxes must increase with $\rho$.

\begin{proposition}[Nesting]\label{prop:nesting}
For $\alpha > 0$ and $0 < \beta < 1/3$, the family
$\{\mathbf{E}_\rho\}_{0 \le \rho \le 1}$ is strictly nested: each shell encloses the
preceding one, and the construction extends continuously to the center, where the
shell at $\rho=0$ degenerates to a point.
\end{proposition}
\begin{proof}
Differentiate: $\frac{ds_x}{d\rho} = 1 + 3\alpha\rho^2 > 0$,
$\frac{ds_y}{d\rho} = 1 > 0$, and $\frac{ds_z}{d\rho} = 1 - 3\beta\rho^2 > 0$, the
final inequality following from $\beta < 1/3$. Strictly increasing semiaxes imply each
shell is interior to the next.
\end{proof}

Near the center the three axes are approximately equal; near the boundary their ratios
differ. We have a family of increasingly elliptical shells.

However, this is only a geometric model. It does not establish that light propagates
along those shells, or that the shells are surfaces of equal optical travel time. To
determine that, we need optics---and a seventeenth-century mathematician enters the
data-structure design meeting.

\section{Fermat's principle as an address generator}\label{sec:fermat}

Fermat's principle states that physical rays follow paths for which optical travel time
is stationary under suitable variations~\cite{fermat1662,bornwolf1999}. For a
refractive-index field $n(\xv)$, the travel-time functional is
\begin{equation}
\mathcal{T}[\gamma] = \frac1c \int_\gamma n(\xv)\, d\ell,
\end{equation}
where $\ell$ is physical arc length. A physical ray satisfies
$\delta\mathcal{T} = 0$; ``stationary'' is more precise than ``minimum,'' since a ray
path need not be a strict global minimum.

Why should a memory engineer care? Because the device is not primarily interested in
where a measurement sits in Euclidean space; it is interested in the phase and arrival
time of optical information. Two spatially nearby measurements can have different
propagation delays; two spatially distant measurements can share the same optical
phase or travel-time surface.

Introduce a phase-travel-time field $T_p(\xv)$. For isotropic geometric optics at a
fixed frequency, away from singularities, it satisfies the eikonal equation:
\begin{equation}
\boxed{\left|\nabla T_p(\xv)\right|^2 = \frac{n_p^2(\xv,\omega_0)}{c^2}.}
\end{equation}
The level sets $\Sigma_\tau = \{\xv : T_p(\xv) = \tau\}$ are equal-phase wavefronts at
the chosen frequency. The memory problem now changes character: instead of asking which
Cartesian cells belong inside an ellipsoid, we ask which optical states belong to each
propagation-related layer. If the actual optical medium produces increasingly
elliptical wavefronts, the resulting memory organization can reflect that deformation.

\begin{remark}
The conditional matters: progressive ellipticity must come from the optical model or
measured calibration data. Fermat's principle does not guarantee it. Using optical
propagation geometry as an organizing principle, rather than treating the ellipsoidal
boundary as a Boolean mask, is our first architectural improvement---and it creates a
new problem: what happens when two rays arrive at the same place?
\end{remark}

\subsection{The Cauchy problem, or why rays refuse to behave like array indices}

Specify an initial wavefront with its initial phase and propagation directions; this is
a Cauchy problem for subsequent propagation. The eikonal equation belongs to the
Hamilton--Jacobi family~\cite{arnold1989,gevay2011}; its characteristic curves derive
from a Hamiltonian. Fix the
optical frequency and suppress it from the notation. Define
\begin{equation}
H(\xv,\pv) = \frac12 \left( |\pv|^2 - \frac{n_p^2(\xv)}{c^2} \right).
\end{equation}
Optical rays lie on $H = 0$ and evolve by Hamilton's equations:
\begin{align}
\frac{d\xv}{ds} &= \frac{\partial H}{\partial \pv} = \pv,\\
\frac{d\pv}{ds} &= -\frac{\partial H}{\partial \xv} = \frac{n_p(\xv)}{c^2} \nabla n_p(\xv),
\end{align}
where $s$ is a characteristic parameter, not physical time. On $H=0$,
$|\pv| = n_p/c$, hence $\left|d\xv/ds\right| = n_p/c$ and $ds = (c/n_p)\,d\ell$, giving
\begin{equation}
\boxed{s - s_0 = c \int_{\ell_0}^{\ell} \frac{d\ell'}{n_p(\ell')}.}
\end{equation}
This parameter is neither physical time nor optical travel time
$T_p = \int \frac{n_p}{c}\, d\ell$; its normalization changes traversal speed but not
the geometric ray trajectory.

Given an initial wavefront parameterized by $(u,v)$, the ray family $\xv(s,u,v)$
locally provides three coordinates for three-dimensional space. That looks promising.
Unfortunately, nature has invented caustics.

\subsection{The caustic problem}

When sunlight passes through a glass of water, bright curves appear on the table
because neighboring rays concentrate. Mathematically, the map
$(s,u,v) \mapsto \xv(s,u,v)$ may cease to be locally invertible. Define the Jacobian
\begin{equation}
J = \det \left[ \frac{\partial \xv}{\partial s},\;
\frac{\partial \xv}{\partial u},\; \frac{\partial \xv}{\partial v} \right].
\end{equation}

\begin{proposition}[Caustic signature]\label{prop:caustic}
The spatial ray-coordinate map loses local invertibility exactly where $J = 0$.
\end{proposition}
\begin{proof}
Immediate from the inverse function theorem: local invertibility of the parametrization
requires a nonsingular Jacobian matrix.
\end{proof}

The key observation: the rays themselves need not cease to exist. What fails is our
attempt to identify a ray using spatial position alone. Suppose two rays reach the same
spatial point, $\xv_1 = \xv_2$; they can still have different momenta,
$\pv_1 \ne \pv_2$. A spatial-only address merges them. An optical state-space address
can keep them separate. This is why the two rotational degrees of freedom in the
original specification are potentially far more important than they first appear: they
can encode the propagation direction that distinguishes ray branches.

\section{The five-dimensional escape hatch}\label{sec:phase-space}

In ordinary Hamiltonian mechanics, a particle in three-dimensional space has a
six-dimensional phase space $(x,y,z,p_x,p_y,p_z)$. In our isotropic fixed-frequency
optical model, the Hamiltonian constraint determines the magnitude of optical momentum,
\begin{equation}
|\pv| = \frac{n_p(\xv,\omega_0)}{c},
\end{equation}
leaving only its direction free. A direction in three dimensions requires two angular
coordinates:
\begin{equation}
\pv = \frac{n_p(\xv,\omega_0)}{c}
\begin{pmatrix} \sin\theta\cos\phi \\ \sin\theta\sin\phi \\ \cos\theta \end{pmatrix}.
\end{equation}
We arrive at the boxed state tuple
\begin{equation}
\boxed{(x,y,z,\theta,\phi)}
\end{equation}
---three spatial dimensions, two rotational dimensions---exactly the five degrees of
freedom from the assignment, under the interpretation that the rotations are
propagation directions. This is the most natural reading for fixed-frequency
propagation in an isotropic medium, but not the only physical interpretation: a
polarization-frame orientation or an aperture sampling scheme could introduce different
angular variables. Our construction specifically uses propagation direction because it
distinguishes intersecting ray branches.

The natural mathematical object is the fixed-frequency constraint manifold
\begin{equation}
\Mopt = \left\{ (\xv,\pv) : |\pv| = \frac{n_p(\xv,\omega_0)}{c} \right\}.
\end{equation}

\begin{proposition}\label{prop:manifold}
For a smooth positive refractive index $n_p$, $\Mopt$ is a five-dimensional manifold,
and the projection $\pi:\Mopt \to \R^3$ may fail to distinguish distinct states: two
rays can occupy the same spatial point but remain distinct points of $\Mopt$ because
their directions differ.
\end{proposition}
\begin{proof}
$F(\xv,\pv) = |\pv|^2 - n_p(\xv,\omega_0)^2/c^2$ satisfies
$\partial_{\pv} F = 2\pv \neq 0$ on $F=0$ (since $n_p>0$ forces $|\pv|>0$), so the
regular level set theorem applies in the six-dimensional ambient space, yielding
dimension $5$. The failure of $\pi$ to be injective is witnessed by any pair
$(\xv,\pv_1)\neq(\xv,\pv_2)$ with $|\pv_i| = n_p(\xv,\omega_0)/c$ and $\pv_1 \neq\pv_2$,
which exists whenever $\xv$ admits two distinct propagation directions.
\end{proof}

\begin{remark}[Why multiple charts]\label{rem:charts}
A five-dimensional coordinate chart is not globally nonsingular: angular coordinates
have poles and seams, and complicated geometry may require multiple charts. Nor is
every ray branch uniquely identified by quantized position and direction: distinct
samples can occupy the same discrete cell, requiring additional sample sequence or
branch identity. Nevertheless the correct basic representation has emerged: store
optical states on the constrained phase-space manifold, rather than storing spatial
cells and attaching orientation as an afterthought.
\end{remark}

\section{Constructing the contiguous address space}\label{sec:addressing}

We want a single memory arena $M_0, M_1, \ldots, M_{N-1}$---no giant five-dimensional
rectangular allocation. Instead, construct a finite partition of the admissible optical
state manifold into cells $C_0, C_1, \ldots, C_{N-1}$, each representing a small region
of position--direction space. The encoder is a map
\begin{equation}
E : \mathcal{M}_{\mathrm{discrete}} \rightarrow \{0,\ldots,N-1\}.
\end{equation}
We require \emph{uniqueness} ($E(q_1)=E(q_2) \implies q_1=q_2$ for canonical discrete
cell identities), \emph{coverage} (every admitted cell has an address), \emph{stability}
(the map remains well-defined near caustics and deformed boundaries), and, where
practical, \emph{geometric locality} (optically nearby states occupy nearby addresses).
The last requirement is an optimization objective, not a consequence of compact
indexing.

\subsection{Multiple optical charts}

A single global chart is fragile. Divide the manifold into local patches and assign each
physical cell a canonical chart owner. For each chart $k$, store its base offset $B_k$,
cell count $N_k$, calibration version, and mapping from canonical cell identities to
local ranks. Define
\begin{equation}
\boxed{E(k,r) = B_k + r, \qquad 0 \le r < N_k,}
\end{equation}
with consecutive intervals $B_0 = 0$, $B_{k+1} = B_k + N_k$. This produces contiguous
storage without gaps. Charts may overlap geometrically, but their canonical ownership
must not, or the same logical cell could acquire multiple addresses.

There is also a crucial distinction between unique cells and unique continuous states:
finite digital storage identifies quantized cells, not every real-valued point inside
them, and the quantization error must be bounded relative to the measurement
requirements.

\subsection{The address budget does not survive refinement unchanged}

A caustic warns that the spatial projection becomes singular; it does not automatically
mean every nearby cell must be refined. Refinement is justified where an error
estimator, ray-branch separation criterion, or measurement sensitivity requires it. And
refinement is not free: with a fixed address budget, allocating more cells near a
difficult region reduces the budget available elsewhere. Two failure modes follow.
First, smooth regions become under-resolved, potentially corrupting timing analysis
where the system was expected to behave quietly. Second, allocation becomes
calibration-dependent: two runs with different optical configurations may generate
different chart tables, changing what a given address means---a reproducibility hazard
disguised as a memory optimization.

The four-day compromise: fix the global cell budget, allocate it by a documented rule,
freeze the resulting chart table during acquisition, and record its version with every
measurement batch. Changing the allocation rule is a change to the instrument, not
merely a software optimization.

\section{Why Ada suddenly makes sense}\label{sec:ada}

At this point a conventional programmer might say, ``Fine, let's put everything in a
hash table.'' The senior infrastructure engineer asks more unpleasant questions: can you
prove the encoder never produces an invalid address? Can you prove two distinct
canonical cells never alias? What happens when an interrupt arrives during a
geometry-table update?

Now Ada becomes interesting. Ada offers strong typing, constrained ranges, and
contracts~\cite{barnes2012}; SPARK, its formally analyzable subset, can establish
selected program properties~\cite{spark2014,mccormick2015}. For precision instrumentation this matters
because a silent indexing error can produce a plausible but physically meaningless
measurement. We do not merely want the program to avoid crashing; we want it to avoid
silently turning one optical state into another.

\subsection{The encoder's proof obligations}

Assume an offline geometry preprocessor has constructed canonical charts and assigned
local ranks. The essential range obligation is easy to understand. We need
$B_k + r < N$ for every valid rank $r$. Given $B_k + N_k \le N$, we obtain
$B_k + r < B_k + N_k \le N$, so the address is in bounds.

\begin{obligation}[Range bound]
For all valid chart identifiers $k$ and ranks $0 \le r < N_k$:
$B_k + r < N$.
\end{obligation}

But range safety is not injectivity.

\begin{theorem}[Injectivity]\label{thm:injective}
If the chart intervals $[B_k, B_k+N_k)$ are pairwise disjoint, then $E(k,r) = B_k + r$
is injective on $\{(k,r) : 0 \le r < N_k\}$; if the intervals additionally cover the
complete address range consecutively, $E$ is a bijection onto that range.
\end{theorem}
\begin{proof}
Suppose $E(i,r) = E(j,s)$. If $i \ne j$, disjoint intervals make equality impossible,
so $i=j$; then $B_i + r = B_i + s$ implies $r=s$. Surjectivity follows from the
consecutive-covering hypothesis by induction over the interval sequence.
\end{proof}

\begin{remark}
These are two related but distinct proof obligations. The range bound concerns a single
call to \texttt{Encode}; injectivity compares arbitrary pairs of valid inputs and is a
\emph{relational} property requiring its own proof, whether by a mathematical argument
or a SPARK ghost lemma. The released specification (Appendix~\ref{app:ada}) supplies
the invariant structure, and its GNATprove transcript discharges all 47 obligations
with no unproved checks: range safety, initialization, and the functional contracts of
\texttt{Encode}, \texttt{Decode}, and the ghost injectivity lemma. This machine-checked
result applies to the released chart table and contract structure; it is not claimed
for arbitrary chart tables or for the geometry preprocessor, and the general statement
remains the mathematical Theorem~\ref{thm:injective}.

There is also a deeper missing proof: nothing in this encoder establishes that the
geometry preprocessor constructed the correct optical partition. The solver establishes
what the cells mean; the verification-ready encoder establishes where they live; the
acquisition pipeline establishes which measurements belong to them. Keeping these
responsibilities separate is what makes the system auditable. Formal verification is
not a property of the author's intentions; it is a claim that must survive actual proof
obligations---which is why the transcript is released alongside the specification.
\end{remark}

\section{The IIT-JEE-style mathematical hammer}\label{sec:hammer}

We temporarily forget the terminology and attack the central ideas with calculus,
vectors, and algebra.

\subsection{Deriving Snell's law from Fermat's principle}

Light travels from point $A$ in a medium of refractive index $n_1$ to point $B$ in a
medium of index $n_2$, through a horizontal interface. Let $h_1,h_2$ be the perpendicular
distances, $L$ the horizontal separation, and $x$ the crossing coordinate. Travel time:
\begin{equation}
T(x) = \frac{n_1}{c}\sqrt{x^2+h_1^2} + \frac{n_2}{c}\sqrt{(L-x)^2+h_2^2}.
\end{equation}
Differentiating and setting to zero:
\begin{equation}
\frac{n_1 x}{\sqrt{x^2+h_1^2}} = \frac{n_2(L-x)}{\sqrt{(L-x)^2+h_2^2}},
\end{equation}
and recognizing $\sin\theta_1 = \frac{x}{\sqrt{x^2+h_1^2}}$,
$\sin\theta_2 = \frac{L-x}{\sqrt{(L-x)^2+h_2^2}}$:
\begin{equation}
\boxed{n_1 \sin\theta_1 = n_2 \sin\theta_2.}
\end{equation}
Snell's law appears directly from stationary travel time. The lesson: physically
meaningful trajectories arise from the optical variational principle, and a programmer's
coordinate grid has no reason to respect those trajectories automatically.

\subsection{Why a caustic destroys spatial invertibility}

Consider the two-dimensional toy family of rays
\begin{equation}
x(t,u) = u(1-t), \qquad y(t,u) = t,
\end{equation}
where $u$ labels the initial ray and $t$ parameterizes its evolution. At $t=1$,
$x(1,u)=0$ for every $u$: all rays meet at $(0,1)$. The Jacobian is
\begin{equation}
J = \det \begin{pmatrix} \partial_t x & \partial_u x \\ \partial_t y & \partial_u y \end{pmatrix}
= \det \begin{pmatrix} -u & 1-t \\ 1 & 0 \end{pmatrix} = t-1,
\qquad \boxed{J = 0 \text{ at } t=1.}
\end{equation}
The inverse mapping fails precisely at the focus: away from $t=1$,
$u = \frac{x}{1-y}$, whose denominator vanishes at $y=1$. This is a mathematical failure
of the coordinates, not a floating-point accident. Now examine the ray tangent
$\frac{\partial(x,y)}{\partial t} = (-u,1)$: at the focus, rays with different $u$ have
different directions. Their spatial positions coincide, but their lifted
position--direction states remain distinguishable---which is why storing orientation
alongside position preserves information that a spatial address loses.

\subsection{Quantifying the danger of a progressively elliptical boundary}

Return to the parametrization of \S\ref{sec:domain}:
$x = \rho(1+\alpha\rho^2)u_x$, $y = \rho u_y$, $z = \rho(1-\beta\rho^2)u_z$. Along the
positive $x$-axis, $x(\rho) = \rho + \alpha\rho^3$ and
$\frac{dx}{d\rho} = 1 + 3\alpha\rho^2$, approaching $1$ at the center and
$\left.\frac{dx}{d\rho}\right|_{\rho=1} = 1+3\alpha$ at the boundary; for
$\alpha=0.2$ this is $1.6$. A fixed radial increment produces approximately $60\%$ more
physical displacement near the outer boundary along this axis than near the center.
Along the $z$-axis, $z(\rho) = \rho - \beta\rho^3$ and
$\frac{dz}{d\rho} = 1 - 3\beta\rho^2$; for $\beta = 0.2$ the boundary value is $0.4$.
Uniform increments in $\rho$ are not uniform physical distances, nor necessarily uniform
optical travel-time increments.

The enclosed volume is
\begin{equation}
V(\rho) = \frac{4\pi}{3}\rho^3 a(\rho)b(\rho)c(\rho),
\end{equation}
and the nesting conditions of Proposition~\ref{prop:nesting} ensure these are actual
nested ellipsoids with increasing enclosed volumes. Differentiating:
\begin{equation}
\frac{dV}{d\rho} = 4\pi\rho^2 abc + \frac{4\pi}{3}\rho^3 (a'bc + ab'c + abc'),
\end{equation}
whose first term resembles ordinary radial growth and whose second term comes from the
changing ellipsoidal shape. A shell-indexing scheme that ignores this contribution
misrepresents the volume distribution---before accounting for refractive index, ray
focusing, or directional sampling. A compact spatial-rank lookup can give every cell a
unique number; it cannot, by itself, ensure those numbers respect the physics.

\section{The femtosecond connection}\label{sec:femto}

One femtosecond is $10^{-15}\,\mathrm{s}$; light travels
$c\Delta t = (3\times10^8)(10^{-15}) = 3\times10^{-7}\,\mathrm{m} = \boxed{300\,\mathrm{nm}}$
during that interval. A propagation-delay error of this scale can contaminate a
measurement intended to resolve femtosecond timing structure.

\subsection{Phase velocity, group velocity, and why one refractive index is not enough}

The eikonal equation of \S\ref{sec:fermat} uses the phase refractive index $n_p$, whose
level sets describe equal phase. For a continuous-wave interferometric measurement,
phase is a natural observable. For a pulse-arrival measurement, however, group delay
becomes important. In a time-independent isotropic medium the dispersion relation is
\begin{equation}
|\mathbf{k}| = \frac{n_p(\xv,\omega)\omega}{c}, \qquad
v_p = \frac{\omega}{|\mathbf{k}|} = \frac{c}{n_p},
\end{equation}
and differentiating,
\begin{equation}
\boxed{v_g = \frac{\partial\omega}{\partial|\mathbf{k}|} = \frac{c}{n_p + \omega\partial_\omega n_p} = \frac{c}{n_g},}
\qquad n_g = n_p + \omega\partial_\omega n_p.
\end{equation}
It would be tempting to replace $n_p$ with $n_g$ inside the eikonal Hamiltonian. That is
generally wrong: the phase dispersion relation determines the rays; the group index is
derived from that dispersion relation, not a replacement for it.

\subsection{The extended Hamiltonian}

Introduce
\begin{equation}
\boxed{\mathcal{H}(\xv,\mathbf{k},\omega) = c^2|\mathbf{k}|^2 - n_p^2(\xv,\omega)\omega^2 = 0,}
\end{equation}
a seven-dimensional description with a codimension-one dispersion constraint. For a
canonical Hamiltonian treatment of propagation and timing, include the time coordinate
$t$ conjugate to $\omega$, giving an eight-dimensional extended phase space before
imposing the constraint. Because the medium is time-independent, frequency is conserved
along rays, $d\omega/d\tau = 0$; at fixed $\omega_0$ the dispersion constraint reduces
to the five-dimensional manifold of \S\ref{sec:phase-space}. Group delay is computed
along the phase-generated ray trajectories: for a narrowband pulse propagating along a
specified ray branch $\gamma$,
\begin{equation}
\boxed{T_g(A,B;\omega_0) = \int_\gamma \frac{n_g(\xv,\omega_0)}{c}\, d\ell,}
\end{equation}
assuming an appropriate isotropic stationary-medium setting and a well-defined ray
branch. The state manifold remains phase-defined; what changes is the functional used to
organize measurements. A group-delay-sensitive memory layout may benefit from grouping
states according to $T_g$ rather than $T_p$. State space and organizing coordinate are
separate architectural decisions.

\subsection{Caustics do not move merely because the observable changes}

Changing the indexing functional from phase delay to group delay does not change the
physical ray family at a fixed frequency, and therefore does not independently move the
physical caustics, which are determined by the spatial projection of the
phase-generated rays. However, a group-delay coordinate may introduce its own
representation difficulties: it may fail to distinguish certain states or become
multivalued when projected onto physical space. Those are coordinate or branch-selection
problems, not new physical caustics. A caustic itself is not necessarily a breakdown of
ray dynamics---rays can remain smooth while their spatial projection becomes
singular. What fails in elementary geometrical optics is often the naive amplitude
approximation, which can diverge at an ideal caustic; uniform asymptotic methods
including Airy-type approximations for folds~\cite{airy1841} and Pearcey-type
approximations for cusps~\cite{pearcey1946,berry1980}
describe caustic neighborhoods without a complete Maxwell simulation. When
wavelength-scale material structure or strong scattering invalidates the approximation,
more complete wave modeling may be necessary.

\subsection{The bandwidth problem}

The time-bandwidth product is constrained by the pulse shape:
$\Delta\omega\,\Delta t \gtrsim O(1)$, with the numerical constant depending on how
duration and spectral width are defined. At optical frequencies, a transform-limited
pulse lasting a few femtoseconds has substantial relative bandwidth; a chirped pulse of
the same temporal duration can have a different spectral width, and chirp alone does not
establish a narrower bandwidth. For a broadband pulse, evaluating $n_g$ only at
$\omega_0$ may be inadequate: higher-order dispersion, spectral phase, and pulse
reshaping can affect the measured arrival time. The four-day prototype therefore assumes
a narrowband regime, uses $T_g$ as a leading-order organizing functional, and explicitly
records that this approximation may not hold for the intended instrument. If the device
sweeps frequency, chart tables may need versioning per band, or frequency may need to
become an additional state coordinate.

The compact lesson: the phase dispersion relation generates the rays; group delay
describes a timing observable along those rays; organizing memory by group delay changes
the indexing geometry, not the underlying ray dynamics. Confusing these roles would mean
building the correct data structure around the wrong physical quantity.

\subsection{The measurement model}\label{sec:measurement}

A simplified timing measurement can be written~\cite{allan1966}
\begin{equation}
t_{\mathrm{obs}} = t_{\mathrm{clock}} + T_g + t_{\mathrm{electronics}} + \epsilon.
\end{equation}
After calibration,
\begin{equation}
r = t_{\mathrm{obs}} - \widehat{T}_g - \widehat{t}_{\mathrm{electronics}},
\end{equation}
and the objective is to investigate whether $r$ contains structure beyond expected noise
and known clock behavior. Our memory architecture can help preserve the information
required for propagation corrections and comparisons; it cannot manufacture measurement
precision. If the detector, digitizer, reference oscillator, or acquisition electronics
cannot resolve the relevant fluctuations, a more sophisticated address mapping cannot
recover information that was never captured.

There is also a subtle question of locality. Two samples may be spatially distant yet
share similar calibrated group delays, making them natural comparison partners for some
analyses. A possible analysis-oriented distance is
\begin{equation}
d^2(q_1,q_2) = (T_1 - T_2)^2 + w_\Omega d_{S^2}^2(\Omega_1,\Omega_2)
+ \frac{\|\xv_1-\xv_2\|^2}{v_g^2},
\end{equation}
where $w_\Omega$ has units of time squared. This is a design metric, not a physical law:
the weights encode which differences matter for a particular analysis, must be
calibrated against the measurement objective, and should not be mistaken for universal
constants of optical geometry.

\section{What can actually be built in four days?}\label{sec:fourdays}

A four-day prototype should establish the architecture's critical invariants and expose
unresolved physics rather than concealing it. The geometry computation need not
initially run on the SoC: it can be performed offline and exported as immutable
calibration tables, separating an expensive numerical solver from a bounded runtime
address lookup. We distinguish three levels of success.

\begin{enumerate}[label=\textbf{Tier \arabic*:}]
\item \textbf{Representation correctness.} Every canonical cell has one valid address,
and decoding reconstructs its identity.
\item \textbf{Geometric correctness.} The cells cover the intended optical state space
with quantified approximation errors, including near difficult boundaries and caustics.
\item \textbf{Measurement usefulness.} The representation improves or preserves
timing-analysis accuracy relative to a suitable baseline.
\end{enumerate}

Day four can demonstrate Tier 1 on a completed test implementation and begin evaluating
Tier 2 on synthetic data; Tier 3 is explicitly outside the four-day scope. Passing Tier
1 does not imply Tier 2, and passing Tier 2 does not imply Tier 3. The chart-table
version must be frozen during an acquisition batch: updating the optical partition
mid-measurement could change what an address means, so a calibration version is part of
the sample's interpretation even if recorded only once per batch. This is the sort of
mundane detail that distinguishes an engineering prototype from a mathematical sketch.

\section{From optical geometry to rate-distortion theory}\label{sec:rate-distortion}

The preceding sections establish that a spatial-only representation can introduce false
timing variation, while a finite-resolution phase-space representation preserves much of
the missing information---but a new problem emerges: how should a limited number of
memory cells be distributed across optical phase space? Uniform spatial sampling is not
necessarily appropriate; nor uniform angular sampling; nor even refining wherever the
ray-coordinate Jacobian becomes small. The correct allocation depends on what the
instrument is trying to measure.

\subsection{The loss function}

Let $Q$ be an optical state drawn from the instrument's operating distribution $\mu$,
and let $T_g(Q)$ denote its group-delay correction. A partition
$\mathcal{P} = \{C_1, \ldots, C_K\}$ assigns each optical state to a discrete cell
$B_{\mathcal{P}}(Q)$. Within each cell, the optimal constant group-delay estimate under
squared error is
$\widehat{T}_g(C_k) = \mathbb{E}[T_g(Q) \mid Q \in C_k]$. The resulting distortion is
\begin{equation}
\boxed{\mathcal{L}(\mathcal{P}) = \mathbb{E}\!\left[ \operatorname{Var}\!\left(T_g(Q) \mid B_{\mathcal{P}}(Q)\right) \right],}
\end{equation}
so for a fixed SoC cell budget $K_{\max}$ the design problem becomes
\begin{equation}
\boxed{\begin{aligned} \mathcal{P}^\star &= \arg\min_{\mathcal{P}} \mathcal{L}(\mathcal{P}), \\ \text{subject to} \quad |\mathcal{P}| &\le K_{\max}. \end{aligned}}
\end{equation}
Additional hardware constraints may require contiguous chart regions, bounded lookup
latency, canonical ownership, minimum spatial resolution, or fixed per-chart
allocations. This is a \emph{task-aware quantization} problem: unlike conventional
geometric meshing, its objective is not merely to approximate the shape of the optical
manifold but to minimize unresolved timing error.

\subsection{Why the variance decomposition matters}

The law of total variance gives
\begin{equation}
\operatorname{Var}(T) = \operatorname{Var}(\mathbb{E}[T \mid B]) + \mathbb{E}[\operatorname{Var}(T \mid B)].
\end{equation}
The first term measures geometric timing variation resolved by the partition; the second
measures timing variation remaining inside cells. Therefore
\begin{equation}
\boxed{\text{Resolved fraction} = 1 - \frac{\mathcal{L}(\mathcal{P})}{\operatorname{Var}(T)}.}
\end{equation}
In one simulation, the empirical decomposition was
$\operatorname{Var}(T) = \operatorname{Var}(\mathbb{E}[T\mid B]) + \mathcal{L}(\mathcal{P})$,
and the partition resolved approximately $92\%$ of the simulated geometric-delay
variance. This percentage uses the empirical total variance from that experiment and
should not be confused with the fraction obtained from earlier finite samples with
different residual estimates. The distinction matters because the loss is an expectation
over the operating distribution, while experiments only estimate that expectation from
finite samples.

\subsection{Calibration is part of the budget}

The optimal cell correction is a conditional expectation, which in practice must be
estimated from calibration data. Suppose cell $k$ contains $m_k$ independent calibration
observations with conditional variance $v_k$, and its correction is estimated by the
sample mean. Then $\operatorname{Var}(\widehat{T}_k \mid C_k) = v_k/m_k$, and under
independent calibration and evaluation data the expected timing mean-squared error
becomes
\begin{equation}
\boxed{\operatorname{MSE} = \sigma^2 + \sum_{k=1}^{K} p_k \left( v_k + \frac{v_k}{m_k} \right),}
\end{equation}
where $p_k = \mathbb{P}(Q \in C_k)$, $v_k = \operatorname{Var}(T_g \mid C_k)$, and
$\sigma^2$ is evaluation detector-noise variance. (If calibration measurements also
contain detector noise, that contribution must be included in the calibration estimator
variance.) This creates a second resource constraint: more cells can reduce within-cell
distortion but also require more calibration observations to estimate their corrections
reliably. With a fixed calibration budget $M$, the design problem becomes a joint
optimization over the partition and its calibration allocation:
\begin{equation}
\boxed{\min_{\mathcal{P},\{m_k\}} \left[ \mathcal{L}(\mathcal{P}) + \sum_k p_k \frac{v_k}{m_k} \right]}
\qquad\text{subject to } |\mathcal{P}| \le K_{\max},\ \sum_k m_k \le M.
\end{equation}
The instrument must budget both memory and calibration effort: a cell that is cheap to
allocate in silicon may be expensive to calibrate.

\begin{proposition}[Optimal calibration allocation]\label{prop:calibration}
If the cell probabilities $p_k$ and conditional variances $v_k$ are known, then the
calibration allocation minimizing the total estimator-variance contribution under a
total sample budget $M$ is, in the continuous relaxation,
\begin{equation}
\boxed{m_k \propto \sqrt{p_k v_k}.}
\end{equation}
\end{proposition}
\begin{proof}
Lagrange multipliers: minimize $\sum_k p_k v_k / m_k$ subject to $\sum_k m_k = M$. The
stationarity condition $-p_k v_k / m_k^2 + \lambda = 0$ gives
$m_k = \sqrt{p_k v_k/\lambda}$; summing and enforcing the constraint fixes $\lambda$.
\end{proof}

So even calibration rays should not necessarily be allocated uniformly---a direct
consequence of the multiplier calculation.

\subsection{The central optimization}

Putting it all together, with the constraints explicit:
\begin{equation}
\boxed{\begin{aligned} \min_{\mathcal{P},\{m_k\}} \quad & \mathbb{E}[\operatorname{Var}(T_g \mid B_{\mathcal{P}})] + \sum_k \frac{p_k v_k}{m_k} \\ \text{subject to} \quad & |\mathcal{P}| \le K_{\max}, \\ & \sum_k m_k \le M, \\ & \mathcal{P} \text{ preserves chart ownership,} \\ & \mathcal{P} \text{ preserves required branch identity,} \end{aligned}}
\end{equation}
with contiguous angular cells as an additional constraint for the dynamic-programming
experiment of \S\ref{sec:experiments}. This formulation makes the central tradeoff
explicit: timing distortion, memory capacity, calibration effort, and optical identity
cannot be optimized independently. For scalar $T_g$ with squared error, the unconstrained
problem is Lloyd--Max quantization~\cite{lloyd1982,max1960}; with contiguity, chart
ownership, and branch-identity constraints, it becomes constrained weighted quantization.
Lloyd--Max provides the foundation, not automatically the global solution: the
unconstrained optimum may merge distinct branches, which is why it cannot be deployed
even when its timing number looks better.

The classical rate-distortion problem asks how much information must be retained to
achieve a desired reconstruction accuracy~\cite{shannon1959,coverthomas}; here the distortion measure
is physically meaningful: unresolved group-delay variation. The hardware imposes a finite
set of addresses, the optical model imposes geometric constraints, and calibration
imposes a statistical cost. The resulting partition is not merely a data structure; it is
part of the measurement system.

\section{Experimental results: when the wrong address becomes a false clock}\label{sec:experiments}

We evaluated five synthetic experiments designed to isolate the consequences of spatial
projection, anisotropic geometry, dispersive propagation, chart-aware memory addressing,
and task-aware allocation. All are evidence type (B) of \S\ref{sec:taxonomy}, except
where an identity is marked (A).

\subsection{Experiment 1: A caustic can manufacture timing structure}

We generated $20{,}000$ rays converging to a common spatial location. Each ray retained
a distinct propagation direction and a direction-dependent group delay. Let $U$ denote
the ray's initial direction parameter, and let the measured arrival time be
\begin{equation}
Y = T_g(U) + \varepsilon,
\end{equation}
where $\varepsilon$ is independent detector noise with $\mathbb{E}[\varepsilon]=0$,
$\operatorname{Var}(\varepsilon)=\sigma^2$. A spatial-only representation merges rays
arriving at the same position; if the reconstruction subtracts a single mean delay, its
residual is
$R_{\mathrm{merged}} = T_g(U) - \mathbb{E}[T_g(U)] + \varepsilon$. By independence,
\begin{equation}
\boxed{\operatorname{Var}(R_{\mathrm{merged}}) = \operatorname{Var}_U(T_g) + \sigma^2.}
\end{equation}
The first term is not clock noise; it is unresolved geometric delay variation. A
phase-space representation retains ray direction, allowing the correct branch-dependent
delay to be subtracted:
\begin{equation}
R_{\mathrm{lifted}} = Y - T_g(U) = \varepsilon,
\qquad \boxed{\operatorname{Var}(R_{\mathrm{lifted}}) = \sigma^2.}
\end{equation}

In the simulation, $\operatorname{std}[T_g(U)] = 2.85\,\mathrm{fs}$,
$\sigma = 0.30\,\mathrm{fs}$, ideal merged residual std $= 2.87\,\mathrm{fs}$, variance
$= 8.21\,\mathrm{fs}^2$. The binned lifted residual with $K=128$ bins had std
$= 0.34\,\mathrm{fs}$, variance $= 0.117\,\mathrm{fs}^2$. Equivalently, under the model,
$\sqrt{\operatorname{Var}_U(T_g) + \sigma^2}/\sigma = 9.55$: the merged residual is
$9.55\times$ the detector-noise std, $91\times$ the variance. Compared to the practical
binned lifted residual, it is $8.34\times$ the std, $69.6\times$ the variance. The
difference between those ratios matters: the first compares to the ideal floor
$\sigma = 0.30\,\mathrm{fs}$; the second compares to the practical floor
$\sqrt{\sigma^2 + \operatorname{Var}(T_g \mid B)} \approx 0.34\,\mathrm{fs}$. The binned
lifted residual is not $\sigma$; it is $\sigma$ plus within-bin geometry
$\operatorname{Var}(T_g \mid B)$. Thus the binned representation resolves
$1 - 0.027/8.12 = 99.7\%$ of the geometric-delay variance in that sample, leaving
$\operatorname{Var}(T_g \mid B) = 0.027\,\mathrm{fs}^2$ unresolved, i.e.\ approximately
$0.3\%$
unresolved.

The experiment demonstrates a specific failure mode: spatial merging can turn
deterministic, branch-dependent propagation delays into apparent stochastic timing
fluctuations. It does not demonstrate a new clock; it demonstrates how a false
clock-like signal can arise.

\paragraph{Figure 1 (numerical summary).} Histograms of spatially merged and
phase-space-corrected timing residuals: the merged distribution is broad and asymmetric
(skew $0.42$, excess kurtosis $-0.85$); the corrected distribution is concentrated near
the detector-noise scale but remains binning-limited (skew $0.05$, excess kurtosis
$0.51$); the Shapiro--Wilk test~\cite{shapiro1965} rejects exact normality on
$n=20{,}000$ at
$p<10^{-6}$. Call it noise-and-binning-limited, not an honest Gaussian.

\subsection{Experiment 2: Uniform shell indices are not uniform physical sampling}

We evaluated the progressively ellipsoidal shell family of \S\ref{sec:domain} with
$\alpha = 0.2$, $\beta = 0.2$. At the outer boundary,
$dx/d\rho = 1.6$ and $dz/d\rho = 0.4$: the selected parameters produce a $4{:}1$ ratio
between the physical displacements generated by equal radial increments along the two
axes. (At $\rho = 0.5$ the ratio is $1.15{:}0.85 \approx 1.35$, recovering $\approx 1$ in
the central nearly spherical region.) Therefore uniformly
spaced $\rho$-bins do not correspond to uniform physical sampling.

This is not an address-collision problem: every cell may still have a unique integer
address. It is a geometric-resolution problem: identical index increments represent
substantially different physical displacements. A chart allocator must therefore account
for geometric distortion, using Jacobian information or an appropriate error-based
refinement criterion. Volume check: $V(1) = \frac{4\pi}{3}(1.2)(1)(0.8) \approx 4.82$
versus sphere $\frac{4\pi}{3} \approx 4.19$; the mean-radius cube estimate
$((1.2+1+0.8)/3)^3 \approx 1$ yields the corresponding comparison, confirming that
shape distortion shifts the enclosed volume distribution beyond naive radial estimates.

\subsection{Experiment 3: Phase delay and group delay are not interchangeable}

We evaluated a synthetic dispersive medium with a prescribed frequency-dependent phase
index, stated for reproducibility:
$n_p(\omega) = 1.401 + 0.006\,(\omega/\omega_0)^2$ (an explicitly anomalous model chosen so
that $\omega\,\partial_\omega n_p > 0$),
$\omega_0 = 2\pi \times 405\,\mathrm{THz}$,
path length $L = 1\,\mathrm{m}$. Then
$n_p(\omega_0) = 1.407$ exactly, and $\omega_0\,\partial_\omega n_p(\omega_0) = 2(0.006) =
0.012$, giving $n_g = n_p + \omega\,\partial_\omega n_p = 1.419$. For a homogeneous path,
$T_p = n_p L / c$ and $T_g = n_g L / c$; using phase delay to correct a group-delay
measurement introduces $(n_g - n_p) L / c \approx 4.0\times10^{-14}\,\mathrm{s} =
40\,\mathrm{fs}$. The synthetic configuration thus produced a $\approx 0.85\%$ relative
timing discrepancy per fixed path length when defined as $(n_g-n_p)/n_g$, and
$\approx 0.86\%$ when defined as $(n_g-n_p)/n_p$; the denominator is stated so the
number is reproducible. This is not a numerical indexing error; it is a mismatch between
the modeled observable and the measured one. The result reinforces the distinction
between the phase dispersion relation, which generates the rays, and group delay, which
supplies the relevant timing correction.

\subsection{Experiment 4: Canonical chart ownership prevents address aliasing}

We tested the article's three-chart address table with
$(B_0, N_0) = (0, 512)$, $(B_1, N_1) = (512, 256)$, $(B_2, N_2) = (768, 256)$; the arena
contains $N = 1024$ addresses. For each valid chart--rank pair we computed
$E(k,r) = B_k + r$. Exhaustive enumeration produced $1{,}024$ addresses, all unique,
covering the complete integer interval $[0, 1024)$. Therefore, for this concrete table,
the encoder is a bijection between valid chart--rank pairs and arena addresses,
instantiating Theorem~\ref{thm:injective} for a finite case.

The experiment also illustrates why the chart identifier is essential: if local ranks
are used as global addresses, cells belonging to different charts alias---in this table,
ignoring the chart id yields $384$ duplicate ranks. This exhaustive test establishes the
bijection for the tested finite table. It does not constitute a general SPARK proof for
arbitrary chart tables, nor does it validate the optical correctness of the chart
partition. Evidence type for this experiment is (B) for the table; (C) is still missing
for the general implementation.

\subsection{Experiment 5: Uniform bins versus task-aware allocation---held-out}

We compared alternative partitions under the same cell budget $K = 64$. The tested delay
function $T_g(\Omega)$ varied slowly over one part of the angular domain and rapidly over
another. A uniform partition allocated equal resolution to both angular regions; a
dynamic-programming partition optimized contiguous angular cells to minimize
$\mathcal{L}(\mathcal{P})$ on training data, subject to the chart-legality constraints
(contiguity, one interval per branch, branch identity recoverable from the address). From
$300{,}000$ rays per seed under the deterministic SHA-256 80/20 split we used the first
$200{,}000$ of the $\approx240{,}000$ available training rays and all
$\approx60{,}000$ held-out rays ($60{,}016\pm223$ across seeds; $60{,}170$ at the
reference seed). Exact held-out improvement for the chart-legal DP partition:
\begin{equation}
\frac{0.000423}{0.000242} = 1.748 \approx \mathbf{1.75\times}
\end{equation}
reduction in within-cell distortion.

\paragraph{What the headline number does and does not mean.}
A single ratio invites over-reading, so we report three controls, all at released
replication cost (12 seeds, $t_{11}=2.201$; \texttt{code/scaling.py}).

\emph{First, the improvement is a property of the budget as much as of the method.}
\begin{center}
\begin{tabular}{rrrrr}
\toprule
$K$ (cells/branch) & uniform $\mathrm{fs}^2$ & DP $\mathrm{fs}^2$ & ratio & 95\% CI \\
\midrule
$16$  ($8$)   & $6.406\times10^{-3}$ & $2.658\times10^{-3}$ & $2.410$ & $\pm 0.011$ \\
$32$  ($16$)  & $1.677\times10^{-3}$ & $8.195\times10^{-4}$ & $2.046$ & $\pm 0.007$ \\
$64$  ($32$)  & $4.231\times10^{-4}$ & $2.421\times10^{-4}$ & $1.748$ & $\pm 0.009$ \\
$128$ ($64$)  & $1.059\times10^{-4}$ & $6.689\times10^{-5}$  & $1.583$ & $\pm 0.005$ \\
$256$ ($128$) & $2.649\times10^{-5}$ & $1.757\times10^{-5}$  & $1.508$ & $\pm 0.007$ \\
\bottomrule
\end{tabular}
\end{center}
The advantage decays with budget and converges towards the generic value in the next
paragraph. Any single-budget figure is therefore a statement about that budget.

\emph{Second, most of the gain is ordinary optimal nonuniform quantization.} At
$K = 64$ (32 cells per branch) the same code applied to other delay laws gives:
\begin{center}
\begin{tabular}{llrr}
\toprule
delay law & interpretation & ratio & 95\% CI \\
\midrule
$A_0\theta$          & affine in $\Omega$            & $0.997$ & $\pm 0.002$ \\
$A_0\sin 2\theta$    & one smooth regime, no split   & $1.504$ & $\pm 0.006$ \\
$A_0 A\sin4\theta$   & the released two-regime law    & $1.748$ & $\pm 0.009$ \\
$A_0\sin 8\theta$    & one fast regime                & $1.889$ & $\pm 0.004$ \\
\bottomrule
\end{tabular}
\end{center}
The affine row is the control that matters most: uniform bins are optimal when the
observable is affine, so the estimator is not simply favouring itself, and $0.997\times$
is the expected null. The smooth-sinusoid row shows that a delay law with \emph{no}
slow/rapid structure at all already delivers $1.50\times$ at the same budget. The
two-regime structure of our released model therefore adds to a generic effect rather
than causing it, and the honest summary is that this experiment confirms the standard
result --- optimal nonuniform quantization beats uniform binning on any curved observable
[9--11] --- within an architecture that must also carry optical identity.

\emph{Third, the sample is not doing the work it appears to do.} The released delay law is
known in closed form, so the identical dynamic program can be run on an analytic
quadrature grid with zero sampled rays. It reaches $2.4192\times10^{-4}\,\mathrm{fs}^2$
against the sample-estimated $2.4208\times10^{-4}\,\mathrm{fs}^2$, a relative gap of
$0.07\%$, with cell boundaries agreeing to within $1$--$3$ grid points out of $1024$. The
partition is a function of the delay law; the $200{,}000$ rays estimate it and add a
little noise. We therefore describe this procedure as partition optimization
\emph{estimated from samples}, not as learning from data, and note that for a calibrated
synthetic benchmark the analytic route is both cheaper and marginally better.

\paragraph{Chart legality is an address-structure cost, not a moral one.}
The sharpest way to price the chart-legal constraint is to compute what it forbids.
Optimizing the scalar observable alone --- the classical optimal scalar quantizer of
$T_g$ (Max 1960 [10]; Lloyd 1982 [9]; Gray--Neuhoff 1998 [11]), precisely the prior work
this paper cites --- with $32$ cells reaches $4.0019\times10^{-5}\,\mathrm{fs}^2$, a
$10.6\times$ improvement over uniform binning rather than $1.75\times$. That solution is
not incorrect; it is unrepresentable. Its cells are contiguous in $T_g$ but not in
$\Omega$, and scanning $\Omega$ they fragment into $91$ runs, so realizing them requires
$\approx2.8\times$ the chart-address structure of the $32$ contiguous angular cells we are
permitted. The earlier framing of this experiment --- an inadmissible competitor that
gains $1.67\times$ by merging opposite branches --- understated the effect, because it
compared a chart-legal optimum against a weak heuristic rather than against the true
scalar optimum.

This yields the result we consider the paper's sharpest statement about representation
cost. The scalar optimum is dramatically better and the chart-legal optimum is
dramatically worse, and the difference is not a moral judgement but a measurable
requirement for address structure. Consequently the design objective of
\S\ref{sec:rate-distortion} is correctly specified: one cannot price a representation by
scalar reconstruction error alone, because topology, identity, address count,
calibration, and distortion are coupled resources. Theorem~\ref{thm:branchcost} prices
one analytic instance of that coupling; the $91/32$ measurement above prices a geometric
instance of the same coupling, and the two are independent routes to the same conclusion.

The released reference generator reproduces both headline distortions
($4.230\times10^{-4}$ and $2.423\times10^{-4}\,\mathrm{fs}^2$) and the $1.746\pm0.005$
ratio over 30 seeds, byte-identically; the scaling, control, and Lloyd--Max tables above
are regenerated by \texttt{code/scaling.py}.

The result illustrates why a purely geometric refinement heuristic can fail: a small
ray-coordinate Jacobian may indicate a caustic, but it does not by itself quantify the
benefit of adding a memory cell. The relevant quantity is the reduction in expected
distortion achieved by refinement,
\begin{equation}
\boxed{\Delta\mathcal{L}(C) = \mathcal{L}(\mathcal{P}) - \mathcal{L}(\mathcal{P} \text{ with } C \text{ refined}),}
\end{equation}
and for a fixed memory budget, refinement should prioritize cells with the greatest
expected reduction in timing distortion per unit of additional storage, subject to
coverage and hardware constraints.

\paragraph{Repeated-seed statistics and released reproduction artifacts.}
The original Experiment~5 run was not replicated, so no uncertainty was originally
reported. We therefore calibrated and released a reference generator
(\texttt{code/generate.py}) whose two calibration constants are documented: an
amplitude fixing the uniform-bin distortion at $4.23\times10^{-4}\,\mathrm{fs}^2$ and a
heterogeneity parameter fixing the DP ratio at the documented value. Over $30$
independent sampling seeds (300k rays each under a deterministic SHA-256 80/20 split,
from which the first 200k training rays and all $\approx60$k held-out rays are used; exact
DP on a 1024-point boundary grid; branch identity enforced by $K/2$ addresses per
branch) the mean held-out distortions were
\begin{equation}
\overline{D}_{\mathrm{uniform}} = 4.230\times10^{-4} \pm 9.0\times10^{-7},
\qquad
\overline{D}_{\mathrm{DP}} = 2.423\times10^{-4} \pm 4.2\times10^{-7}
\end{equation}
(95\% CI, $t_{29}$), giving a per-seed improvement ratio
$\mathbf{1.746 \pm 0.005}$ (95\% CI). The ratio is stable across seeds; the absolute
distortion scale is set by the calibrated amplitude, and the confidence interval
quantifies split and sampling variability under the released model, not model-selection
uncertainty. All per-seed values are in \texttt{code/results\_exp5\_seeds.csv}.

\paragraph{Figure 2 (numerical summary).} Held-out within-cell distortion at equal cell
budget $K=64$, on the released calibrated model: DP improves $1.75\times$ over uniform
angular binning, while optimizing the scalar observable alone would improve $10.6\times$
but requires $91$ angular runs instead of $32$.

\paragraph{One clarification about scope.} The released delay law is branch-independent
and is not $\pi$-symmetric: $T_g(\theta) \neq T_g(\pi - \theta)$ for all
$\theta \neq \pi/2$. This generator therefore \emph{cannot} exhibit a
$\theta \leftrightarrow \pi-\theta$ branch collision, and any contiguous angular
partition preserves branch identity trivially on it. Branch identity is accordingly a
theoretical architectural constraint throughout this paper
(Sec.~\ref{sec:addressing}, Theorem~\ref{thm:branchcost}), not a quantity Experiment~5 measures; what Experiment~5
measures is chart-legal contiguous quantization of a scalar observable. We removed an
earlier draft claim that a released transform-based partition merged
$\theta = 0.35\pi$ with $\theta = 0.65\pi$ because they share a delay value, as the
released model does not support it.

\subsection{Calibration allocation at large and small M}

Using the chart-legal DP bins of the previous subsection, with $p_k$ and $v_k$ estimated
from training data and the
continuous relaxation of Proposition~\ref{prop:calibration}: the optimal allocation
offers negligible improvement when calibration is abundant but becomes consequential when
samples are scarce, because the calibration term $p_k v_k / m_k$ grows as $1/m_k$ while
the within-cell term stays fixed. At $M = 10^6$, optimal $\{m_k\}$ is nearly uniform;
at $M = 256$, optimal allocations deviate strongly from uniformity, showing the
continuous relaxation breaks down and integer allocation is required. Finite-calibration
simulation for the binned lifted estimator with floor variance
$0.117\,\mathrm{fs}^2$ (calibration-term inflation $v_k/m_k$, $v_k \approx 0.117$):
$m = 1$ gives $+100\%$ (variance $0.234\,\mathrm{fs}^2$), $m = 2$ gives $+50\%$,
$m = 4$ gives $+25\%$, and $m = 16$ gives $+6.25\%$; article.md's $+16.6\%$ corresponds to
$m = 6$. A cell cheap in silicon may be expensive to calibrate.

\section{Discussion: the memory problem was never just about memory}\label{sec:discussion}

We began with five logical dimensions and a contiguous memory requirement. Ordinary
row-major indexing solved the dimensionality problem but ignored the domain's geometry.
A compact spatial-rank lookup solved wasted allocation but did not establish the locality
and sampling properties needed by the optical application. Progressive ellipticity forced
us to examine the geometry itself. Fermat's principle introduced optical travel time; the
eikonal equation produced equal-phase wavefronts; the Cauchy problem and Hamiltonian
characteristics described how those wavefronts propagate; caustics showed why spatial ray
coordinates can lose invertibility; the fixed-frequency optical phase space explained the
natural appearance of three positional and two directional degrees of freedom; a
multi-chart partition provided a route to contiguous, collision-free digital addresses;
Ada and SPARK gave us tools to express and potentially verify the discrete
representation's invariants; and the phase-versus-group-delay distinction forced the
separation between the manifold, generated by the phase dispersion relation, and group
delay, a different observable computed along its rays. The memory layout must choose
which functional it organizes by: for a pulse-arrival observable, group delay is a
natural leading-order choice, but that does not license rebuilding the manifold from the
group index. State space and organizing coordinate are separate decisions.

The atomic-clock application explains why all this matters: the system must not confuse
geometric propagation effects with genuine timing fluctuations. There is one last
mathematical distinction worth appreciating. Hamiltonian evolution preserves symplectic
structure in full phase space; under suitable conditions, Liouville's
theorem~\cite{liouville1838} tells us
that phase-space volume is preserved by the flow. But preservation of full phase-space
volume does not mean spatial volume is preserved: a ray bundle can focus dramatically in
position space while spreading in momentum directions. That is why caustics are
compatible with well-behaved Hamiltonian evolution, and it is why a representation that
keeps only spatial coordinates can become inadequate. The mathematics was quietly telling
us where to put the missing information: not in another arbitrary array---in the
directions of the rays.

\subsection{Interpretation of all five experiments}

The five experiments expose five different failure modes at different layers of the
system; fixing one does not automatically fix the others.

\begin{itemize}[nosep]
\item \textbf{Spatial projection} can manufacture apparent timing fluctuations by merging
physically distinct ray branches (Exp.~1).
\item \textbf{Uniform logical sampling} can misrepresent physical resolution in an
anisotropic domain (Exp.~2).
\item A \textbf{phase-versus-group-delay mismatch} can introduce systematic timing error
even when the data structure itself is correct (Exp.~3).
\item Omitting \textbf{canonical chart ownership} can produce actual memory-address
collisions (Exp.~4).
\item \textbf{Unconstrained task-aware optimization} can improve a single timing number
while destroying branch identity (Exp.~5).
\end{itemize}

The central experimental conclusion is therefore not that a five-dimensional array is
inherently superior to a three-dimensional array. It is that preserving the optical state
variables required by the measurement prevents information loss that a spatial-only
representation cannot recover after projection. The experiments establish this result in
controlled synthetic settings; physical validation, broadband propagation modeling, and
machine-checked verification remain separate obligations.

\section{Five theoretical extensions}\label{sec:extensions}

The five results in this section are deductions from the model of the preceding
sections, with each additional assumption stated explicitly. The first limits the
five-dimensional representation itself.

\subsection{Closed-ray obstruction: even exact position and direction may not determine accumulated delay}

Suppose an admitted ray returns to the same position \emph{and} direction after a
complete circuit, with positive group index along the circuit.

\begin{theorem}[Closed-ray obstruction]\label{thm:closedray}
Accumulated group delay cannot be represented by a single-valued function of the
instantaneous five-dimensional optical state on any domain that includes repeated
circuits.
\end{theorem}
\begin{proof}
Let $q(s)$ denote the position--direction state, with $q(P) = q(0)$ after one circuit of
period $P$. If accumulated delay were a state function $G(q)$, returning to the same
state would imply $G(q(P)) - G(q(0)) = 0$. But propagation around the circuit accumulates
$\Delta T_g = \oint_\gamma \frac{n_g}{c}\,d\ell > 0$. The same difference would have to
be both zero and positive---a contradiction.
\end{proof}

This is a continuous-state obstruction: increasing angular or spatial resolution cannot
fix it. The caustic argument of \S\ref{sec:phase-space} distinguishes rays with
different directions; this result identifies a further ambiguity: two visits can have
identical position and direction but different travel times. The five-dimensional
representation can identify the instantaneous ray state while failing to identify the
measurement's history, so a recirculating instrument needs additional information such
as a lap counter, launch-event identity, or accumulated delay. Single-pass acquisition
avoids this obstruction by selecting one traversal from an entry surface.

\subsection{Hidden branches can manufacture an entire clock-like noise spectrum}

The paper proves that unresolved geometry adds variance; a stronger result concerns its
\emph{temporal structure}. Suppose two branches share one stored address with delays
$T_g(Z_t) = T_0 + aZ_t$, $Z_t \in \{-1,+1\}$, where the branch sequence is a stationary
symmetric Markov chain with
$\Pr(Z_{t+1} = Z_t) = \frac{1+\rho}{2}$, $0 < \rho < 1$. Assume a perfectly stable
clock and independent white detector noise $\varepsilon_t$ of variance $\sigma^2$.
After subtracting the merged-cell correction $T_0$, the residual is
$R_t = aZ_t + \varepsilon_t$.

\begin{theorem}[Clock-like spectrum]\label{thm:spectrum}
The residual has autocovariance and discrete-frequency power spectrum
\begin{equation}
\boxed{
\operatorname{Cov}(R_t,R_{t+h}) = a^2\rho^{|h|}+\sigma^2\mathbf{1}_{\{h=0\}},
\qquad
S_R(\Omega) = \sigma^2 + a^2\frac{1-\rho^2}{1-2\rho\cos\Omega+\rho^2}.
}
\end{equation}
\end{theorem}
\begin{proof}
The transition rule gives $\mathbb{E}[Z_{t+h} \mid Z_t] = \rho^h Z_t$. Since
$\mathbb{E}[Z_t] = 0$ and $Z_t^2 = 1$, the covariance follows directly; Fourier
transforming the geometric sequence gives the spectrum.
\end{proof}

At zero frequency the geometric contribution is
$S_{\mathrm{geom}}(0) = a^2\frac{1+\rho}{1-\rho}$: as branch persistence approaches
one, substantial low-frequency power appears even though the clock is perfect. There is
also an identifiability consequence: without additional restrictions or branch
observations, the explanations (i)~stable clock with fluctuating unresolved branch delay
$aZ_t$ and (ii)~constant propagation delay with clock fluctuation $aZ_t$ produce exactly
the same recorded residual process. Temporal correlation alone cannot distinguish hidden
optical geometry from clock fluctuations; validation should therefore test
autocorrelation and spectra after branch correction, alongside the residual variance
already considered, using the metrological machinery of Allan
statistics~\cite{allan1966}.

\subsection{Preserving branch identity has an exact distortion cost}

Consider a controlled model: $J$ equally probable branches, delay $T$ uniformly
distributed over an interval of width $W$ on every branch, branch identity independent
of delay, $K$ total addresses with $K$ divisible by $J$, one reconstructed delay per
address, and branch identity recoverable from the address. For a uniform scalar, the
optimal $k$-cell squared-error distortion is $W^2/(12k^2)$, obtained by equal intervals
and midpoint reconstruction~\cite{lloyd1982,gray1998}.

\begin{theorem}[Branch-preservation cost]\label{thm:branchcost}
Subject only to exact branch preservation,
$D^\star_{\mathrm{preserve}}(K) = \frac{J^2 W^2}{12 K^2}$; if branch identity may be
discarded, $D^\star_{\mathrm{discard}}(K) = \frac{W^2}{12K^2}$; consequently
\begin{equation}
\boxed{\frac{D^\star_{\mathrm{preserve}}}{D^\star_{\mathrm{discard}}} = J^2.}
\end{equation}
\end{theorem}
\begin{proof}
Exact branch recovery requires disjoint address sets. Let branch $b$ receive $K_b$
addresses with $\sum_{b=1}^J K_b = K$; its conditional distortion is at least
$W^2/(12K_b^2)$, so
$D \ge \frac{W^2}{12J}\sum_{b=1}^J \frac{1}{K_b^2}$. Convexity of $x^{-2}$ makes equal
allocation $K_b = K/J$ optimal; substitution proves the result, and equal delay
intervals attain it.
\end{proof}

In this model, preserving four independent branch identities costs a factor of $16$ in
MSE---or $4$ in RMS error---at the same total address capacity. An inadmissible method
can therefore win a timing benchmark by a large factor purely because it discards
required information: the counterexample of Experiment~5 has an exact theoretical
counterpart. The budget must count the combined identity; storing branch labels
separately consumes additional bits even when correction tables can be shared. This
places the branch-constrained design inside the functional-quantization
conversation~\cite{misra2011}, where quantizers minimize the error of a computed
function rather than a raw reconstruction, and into the broader quantization
theory~\cite{shannon1959,gray1998}.

\subsection{Refinement has a calibration threshold---and useful resolution grows only as \texorpdfstring{$M^{1/3}$}{M\^{}(1/3)}}

For partition $\mathcal{P}$ define
$p_k = \Pr(Q \in C_k)$ and $v_k = \operatorname{Var}(T_g \mid C_k)$; let each correction
be estimated from $m_k$ independent calibration observations, each carrying additional
independent noise of variance $\nu^2$, and write
$D(\mathcal{P}) = \sum_k p_k v_k$ and
$S(\mathcal{P}) = \sum_k \sqrt{p_k(v_k + \nu^2)}$.

\begin{theorem}[Calibration threshold]\label{thm:threshold}
Under a total calibration budget $\sum_k m_k = M$, the continuous allocation optimum is
\begin{equation}
\boxed{
R^\star(\mathcal{P}) = \sigma^2 + D(\mathcal{P}) + \frac{S(\mathcal{P})^2}{M},
}
\end{equation}
where $\sigma^2$ is evaluation detector-noise variance.
\end{theorem}
\begin{proof}
Independent calibration and evaluation give
$R = \sigma^2 + D + \sum_k \frac{p_k(v_k+\nu^2)}{m_k}$. Cauchy--Schwarz yields
$\bigl(\sum_k \sqrt{p_k(v_k+\nu^2)}\bigr)^2 \le
\bigl(\sum_k \frac{p_k(v_k+\nu^2)}{m_k}\bigr)\bigl(\sum_k m_k\bigr)$, with equality for
$m_k = M\,\frac{\sqrt{p_k(v_k+\nu^2)}}{S(\mathcal{P})}$, the continuous relaxation of
Proposition~\ref{prop:calibration} extended by $\nu^2$; actual sample counts require
integer allocation.
\end{proof}

Consequently, a refinement $\mathcal{P} \to \mathcal{P}'$ improves the relaxed total
error \emph{if and only if}
\begin{equation}
\boxed{
D(\mathcal{P}) - D(\mathcal{P}') > \frac{S(\mathcal{P}')^2 - S(\mathcal{P})^2}{M}.
}
\end{equation}
Reducing geometric distortion alone is insufficient. For an exact counterexample, let
the true delay be constant everywhere: with $K$ equally probable cells, separate
sample-mean calibration gives $R - \sigma^2 = K\nu^2/M$---every additional cell increases
error while providing no geometric benefit.

A useful scaling law follows for an equal-probability design family with within-cell
variance $A/K^2$. Equal calibration allocation gives
$R(K) - \sigma^2 = \frac{A}{K^2} + \frac{A}{MK} + \frac{\nu^2 K}{M}$; for $K \ll M$ and
$\nu^2 > 0$, balancing the dominant terms gives
\begin{equation}
\boxed{
K_{\mathrm{opt}} \approx \left(\frac{2AM}{\nu^2}\right)^{1/3},
}
\end{equation}
subject to memory and sampling feasibility. Doubling useful resolution therefore
requires approximately eight times the calibration budget; the $M^{-2/3}$ error scaling
connects to classical nonparametric estimation rates~\cite{stone1982}. The refinement
rule of \S\ref{sec:addressing} should thus optimize expected \emph{measurement-error}
reduction after calibration, rather than geometric distortion reduction alone.

\subsection{Ray-path sensitivity contributes to dispersion through an inverse Hessian}

This strengthens the phase-versus-group-delay discussion of \S\ref{sec:femto}.
Parameterize admissible fixed-endpoint paths by a finite vector $z$ with redundant
coordinates removed, and define the accumulated phase
\begin{equation}
\Phi(z,\omega) = \frac{\omega}{c}\int_{\gamma_z} n_p(\xv,\omega)\,d\ell.
\end{equation}
A smooth physical branch satisfies $\nabla_z \Phi(z_\star(\omega),\omega) = 0$; assume
the path Hessian $H = \Phi_{zz}$ is invertible. For the stationary-path phase
$\phi(\omega) = \Phi(z_\star(\omega),\omega)$, group delay is $\phi'(\omega)$ and
group-delay dispersion is $\phi''(\omega)$~\cite{bonnans2000}.

\begin{theorem}[Dispersion sensitivity]\label{thm:dispersion}
On the stationary branch,
\begin{equation}
\boxed{
T_g = \phi'(\omega) = \Phi_\omega = \int_{\gamma_{z_\star}} \frac{n_g}{c}\,d\ell,
\qquad
\frac{dT_g}{d\omega} = \Phi_{\omega\omega} - \Phi_{\omega z} H^{-1} \Phi_{z\omega}.
}
\end{equation}
\end{theorem}
\begin{proof}
Differentiating stationarity: $H z_\star'(\omega) + \Phi_{z\omega} = 0$, so
$z_\star' = -H^{-1}\Phi_{z\omega}$. The first derivative of the stationary-path phase is
$\phi' = \Phi_\omega + \Phi_z z_\star' = \Phi_\omega$ because $\Phi_z = 0$ on the
stationary branch. Differentiating once more and substituting $z_\star'$ proves the
dispersion formula.
\end{proof}

The important asymmetry: path variation cancels from the \emph{first} phase derivative
but returns in the \emph{second}. If $H$ has a small eigenvalue $\lambda$ with unit
eigenvector $e$, the path contribution includes $-|e^\top \Phi_{z\omega}|^2/\lambda$,
which can become large when frequency couples to a weakly constrained path deformation.
The formula applies while the branch remains nondegenerate; at singularity, a different
local treatment is needed. Assessing bandwidth using material dispersion alone can
therefore miss sensitivity caused by frequency-dependent ray bending. For a smooth
branch over $|\omega - \omega_0| \le h$, the mean-value theorem supplies the practical
bound
\begin{equation}
\boxed{
\left|T_g(\omega) - T_g(\omega_0)\right| \le
h \sup_{\text{band}} \left|\Phi_{\omega\omega} - \Phi_{\omega z} H^{-1} \Phi_{z\omega}\right|,
}
\end{equation}
and requiring this to stay below the timing tolerance gives a concrete criterion for
when one frequency-specific calibration table is adequate and when the frequency
partition of \S\ref{sec:femto} needs refinement.

\subsection{What the extensions change in the design}

\begin{itemize}[nosep]
\item \textbf{Extension 1:} recirculating instruments need history variables (lap counter,
launch identity, or accumulated delay) beyond the five-dimensional state.
\item \textbf{Extension 2:} validation must test autocorrelation and spectra after branch
correction, not only residual variance; clock-like spectra can be geometric artifacts.
\item \textbf{Extension 3:} the benchmark budget must count branch identity exactly;
$J$ branches cost $J^2$ in MSE relative to identity-discarding compression.
\item \textbf{Extension 4:} refinement must clear the calibration threshold
of Theorem~\ref{thm:threshold}; the paper's refinement rule should target
post-calibration error, and constant-delay cells are strictly harmful.
\item \textbf{Extension 5:} bandwidth assessment needs the inverse-Hessian dispersion
sensitivity, not material dispersion alone.
\end{itemize}

\section{Limits and unproved obligations}\label{sec:limits}

By Thursday afternoon, the prototype has produced: a reproducible optical model, a finite
phase-space partition, a contiguous Ada memory arena, an explicit set of address
invariants, and a test harness designed to send rays through a caustic. It knows the
difference between two rays meeting at the same place; it knows that outer shells do not
behave like inner shells; it knows that the phase index and the group index are not
interchangeable; and it knows that a successful memory lookup is not proof of a
successful physical measurement. It has not discovered a new fundamental clock, the SoC
is not fabricated, and nobody has proved that nature contains a subtler clock beneath the
atomic clock.

Three things remain unsettled, and we state them as specific, testable claims rather than
as failures:
\begin{enumerate}[label=(\arabic*)]
\item \textbf{Implementation verification scope.} The released Ada/SPARK
specification is machine-checked by GNATprove (47/47 proof obligations discharged, 0
unproved, for the released three-chart encoder instance and its contract structure;
Appendix \ref{app:ada}), which closes Tier~1 for the released chart table. The chart
table is a three-entry literal, so injectivity holds there by construction; the table
consistency predicate is not used as any subprogram's precondition, the ghost
\texttt{Global\_Rank} entry point is an unimplemented stub, and the exhaustive bijection
test is executed rather than verified. This does not cover the geometry preprocessor,
the interrupt-time chart-table update path, or arbitrary chart tables; those remain
outside the transcript's scope and are stated as such.
\item \textbf{Geometric preprocessing quality.} The quality of the geometric
preprocessing---quantified cell coverage and approximation error near caustics and
deformed boundaries---is a numerical question the four-day schedule does not close
(Tier 2, \S\ref{sec:fourdays}).
\item \textbf{Metrological benefit.} Whether this architecture improves a real
femtosecond measurement over a simpler baseline is a scientific question no data
structure can answer on its own (Tier 3).
\end{enumerate}

These are the criteria on which the prototype must eventually be judged. We close with
the paper's conclusion, restated once more: the optimal representation is not the one
that stores the fewest coordinates or minimizes a single timing error; it is the least
expensive representation that preserves every distinction the instrument may need to
make. The intern started with an array and ended with a constrained
information-theoretic instrument-design problem.

\section*{Acknowledgments}
The tools obscura (a lightweight headless browser) and TinyFish CLI (web search and
fetch) were used during the preparation of this manuscript for source retrieval and
bibliographic verification.

\appendix

\section{Ada/SPARK specification and GNATprove transcript}\label{app:ada}

\subsection{Specification (\texttt{chart\_encoder.ads})}

\begin{verbatim}
--  Chart-aware contiguous memory encoder for discretized optical phase space.
--  Companion artifact to "Chart-Aware Contiguous Memory for Discretized Optical
--  Phase Space" (Sethurathienam Iyer, 2026).
--
--  Design intent (paper Sections 7-8):
--    * E(k, r) = B(k) + r   with consecutive intervals B(0) = 0,
--      B(k+1) = B(k) + N(k).
--    * Range bound        : every address produced is < Arena_Size.
--    * Injectivity        : distinct (chart, rank) pairs yield distinct addresses
--      (relational property, proved via ghost lemma in body).
--    * Coverage/bijectivity follows from consecutiveness of the intervals.
--
--  The chart table below is the concrete table used in Experiment 4:
--  (B0, N0) = (0, 512), (B1, N1) = (512, 256), (B2, N2) = (768, 256).

package Chart_Encoder with SPARK_Mode is

   Chart_Count : constant := 3;
   Arena_Size  : constant := 1024;

   type Chart_Id  is new Natural range 0 .. Chart_Count - 1;
   type Rank      is new Natural range 0 .. Arena_Size - 1;
   type Address   is new Natural range 0 .. Arena_Size - 1;

   --  Chart table (immutable at acquisition time; versioned per batch).
   Base     : constant array (Chart_Id) of Natural := (0, 512, 768);
   Cells    : constant array (Chart_Id) of Natural := (512, 256, 256);

   --  Contract on the raw table: intervals are consecutive and in bounds.
   Table_Is_Consistent : constant Boolean :=
     Base (0) = 0
     and then (for all J in Chart_Id =>
                 Base (J) + Cells (J) <= Arena_Size)
     and then (for all J in Chart_Id'First .. Chart_Id'Last - 1 =>
                 Base (Chart_Id'Succ (J)) = Base (J) + Cells (J));

   --  Total number of cells across all charts.
   Cell_Total : constant := 512 + 256 + 256;  --  = Arena_Size

   type Valid_Rank is new Natural range 0 .. Cell_Total - 1;

   --  The map from linear ranks within a chart to addresses is injective
   --  globally: this predicate expresses that (K, R) corresponds to a
   --  unique global rank, hence a unique address.
   function Global_Rank (K : Chart_Id; R : Valid_Rank) return Natural with
     Pre     => Natural (R) < Cells (K),
     Global    => null,
     Import,
     Ghost;
   --  Encoder: the ONLY address-producing primitive.
   --  The rank must lie within the chart's own cell count; the range bound
   --  then follows from the consecutive-interval table invariant.
   function Encode (K : Chart_Id; R : Valid_Rank) return Address with
     Pre  => Natural (R) < Cells (K)
               and then Base (K) + Natural (R) < Arena_Size,
     Post => Natural (Encode'Result) = Base (K) + Natural (R);

   --  Decoder: reconstructs the canonical cell identity from an address.
   --  Establishes Tier 1 representational correctness: round trip.
   procedure Decode
     (A : Address;
      K : out Chart_Id;
      R : out Valid_Rank) with
     Post => Encode (K, R) = A;

end Chart_Encoder;
\end{verbatim}

\noindent The body adds a ghost lemma \texttt{Injectivity\_Lemma(K1,R1,K2,R2)} whose
precondition includes the pairwise-disjoint consecutive-interval table invariant and
whose postcondition concludes $K_1=K_2 \wedge R_1=R_2$.

\subsection{Transcript (GNATprove 14.1.1; cvc5, Alt-Ergo, Z3)}

\begin{verbatim}
Summary of SPARK analysis
---------------------------------------------------------------------------------
SPARK Analysis results        Total      Flow      Provers   Justified   Unproved
---------------------------------------------------------------------------------
Initialization                    2         2            .           .          .
Run-time Checks                  29         .    29 (CVC5)           .          .
Assertions                        9         .     9 (CVC5)           .          .
Functional Contracts              6         .     6 (CVC5)           .          .
Termination                       1         1            .           .          .
---------------------------------------------------------------------------------
Total                            47    3 (6%)     44 (94%)           .          .

  Chart_Encoder          flow analyzed and proved (4 checks)
  Chart_Encoder.Encode   flow analyzed and proved (13 checks)
  Chart_Encoder.Decode   flow analyzed and proved (6 checks)
  Chart_Encoder.Injectivity_Lemma  analyzed and proved (21 checks)
\end{verbatim}

\noindent The exhaustive runtime test (\texttt{test\_encoder.adb}) additionally
enumerates all 1024 chart--rank pairs, confirms no duplicate addresses, checks that the
union covers $[0,1024)$, and verifies the decode round trip. Both the specification and
the transcript are released with the paper (\texttt{code/chart\_encoder.*},
\texttt{code/gnatprove-transcript.txt}). Scope: the proof applies to this chart table
and contract structure; it does not cover the geometry preprocessor or the
interrupt-time update path (\S\ref{sec:limits}).

\section{Contact geometry of the fixed-frequency manifold}\label{app:contact}



The fixed-frequency optical manifold $\Mopt$ is a regular hypersurface in the cotangent
bundle of three-dimensional space. The ambient cotangent bundle has a canonical
one-form $\lambda = \pv \cdot d\xv$; restrict it to the hypersurface,
$\alpha = \lambda|_{\Mopt}$. For a smooth positive isotropic refractive index, this
defines a \emph{contact form}~\cite{geiges2008}: its kernel is a maximally non-integrable
four-dimensional
plane field on the five-dimensional manifold. The associated Reeb vector field is
tangent to the optical characteristic flow, its trajectories coinciding with the
phase-generated ray trajectories up to reparameterization. The full six-dimensional
phase space is symplectic, while its fixed-frequency optical hypersurface carries a
natural contact structure; the relationship can be understood through symplectization,
the positive radial extension of the contact hypersurface.

This explains why five dimensions arise so naturally: they are not merely the result of
deleting one momentum coordinate from a six-dimensional array; they are the dimension of
a contact manifold whose distinguished flow describes optical rays. None of this is
necessary for the four-day prototype. It suggests a deeper research direction---whether
a discrete optical state representation can preserve useful contact-geometric invariants
with controlled approximation error---which is far more ambitious than whether an Ada
array can store five indices. It also explains why chart-aware addressing is
geometrically natural, even though contact geometry alone does not prescribe a unique or
optimal memory layout.


\section{Phase 2: closed rays, a dynamical zeta, and an obstruction}\label{app:phase2}

Appendix~\ref{app:contact} put a contact form $\alpha = \pv\cdot d\xv$ on the
five-dimensional fixed-frequency manifold, with the Reeb flow recovering the optical
characteristic flow up to reparameterization. Section~\ref{sec:extensions} proved the
closed-ray obstruction (Theorem~\ref{thm:closedray}). This appendix takes the obvious
next step and asks what the \emph{global} closed-ray structure is, in the standard
dynamical-systems sense. It is explicitly speculative: nothing here is required for the
four-day prototype, and the section's purpose is to mark the boundary between what the
geometry already gives and what it does not.

\paragraph{Periods as invariants.} Let $\mathcal{P}$ be the set of primitive closed
Reeb orbits, i.e.\ closed optical rays traversed once. Along $\gamma \in \mathcal{P}$
the Reeb condition $\alpha(R_\alpha) = 1$ identifies Reeb time with the contact action,
\begin{equation}
T^{(p)}_\gamma = \oint_\gamma \alpha
= \oint_\gamma \frac{n_p(\xv,\omega_0)}{c}\,d\ell ,
\end{equation}
which is the phase travel time, and the group-delay observable supplies a second
invariant on the same orbits,
$G_\gamma = \oint_\gamma (n_g/c)\,d\ell$. These are not independent data. With the
positive dispersion ratio $r = n_g/n_p$,
\begin{equation}
G_\gamma = \int_0^{T^{(p)}_\gamma} r\!\left(\phi^t q\right) dt ,
\end{equation}
so group delay is a \emph{positive time change} of the phase period: the group-delay
spectrum is the phase spectrum pushed through $r$, not a fresh object
[evidence type A].

Theorem~\ref{thm:closedray} then reads as a statement about repetition. After $m$
circuits, $q_m = q_0$ while $T_g(m) - T_g(0) = m\,G_\gamma$. The instantaneous
five-dimensional state cannot distinguish $m=1$ from $m=2$, because both are the same
point; the missing history variable \emph{is} the repetition index of a periodic orbit.
That is precisely the primitive-orbit/repetition decomposition that dynamical zeta
functions exist to encode.

\paragraph{A dynamical zeta function.} For a reference time $T_*$, define
\begin{equation}
\zeta_{\mathrm{opt}}(s) = \prod_{\gamma\in\mathcal{P}}
\left(1 - e^{-sG_\gamma/T_*}\right)^{-1},
\qquad
-\frac{\zeta_{\mathrm{opt}}'(s)}{\zeta_{\mathrm{opt}}(s)}
= \sum_{\gamma\in\mathcal{P}}\sum_{m\ge1}\frac{\log N_\gamma}{N_\gamma^{ms}},
\end{equation}
with orbit norm $N_\gamma = \exp(G_\gamma/T_*)$. The Euler product, the repetition
index $m$, and the log-weights therefore assemble themselves exactly as in the Riemann
Euler product. This is why trace formulas resemble the explicit formula, and it is also
why the resemblance is not by itself evidence of anything: every zeta function has an
Euler product over primitive objects with a repetition index. The periodic-orbit measure
$\Theta_{\mathrm{opt}}(t) = \sum_{\gamma}\sum_{m\ge1}G_\gamma\,\delta(t-mG_\gamma)$ has
Laplace transform $-\zeta_{\mathrm{opt}}'/\zeta_{\mathrm{opt}}$, so the zeta is a
packaging of closed-ray delays rather than an extra assumption.

\paragraph{An obstruction.} The cheapest attempt at an arithmetic construction fails,
and this is the one provable result of this appendix. Replace the medium by a finite
optical graph with $E$ edges of fixed optical lengths $\ell_1,\dots,\ell_E$. Each closed
orbit traverses edge $i$ a non-negative integer number of times, so its period lies in
the finite-rank additive module $\Lambda = \mathbb{Z}\ell_1 + \cdots + \mathbb{Z}\ell_E$.
But the prime logarithms are linearly independent over $\mathbb{Q}$: if
$\sum_j q_j\log p_j = 0$, clearing denominators gives $\sum_j m_j\log p_j = 0$ with
$m_j\in\mathbb{Z}$, and exponentiating gives $\prod_j p_j^{m_j} = 1$, so unique
factorization forces every $m_j = 0$. Hence
$\dim_{\mathbb{Q}}\operatorname{span}\{\log p\} = \infty$, and a rank-$E$ module cannot
contain it.

\begin{theorem}[Finite-graph obstruction]\label{thm:finitegraph}
No finite optical graph with additive edge-length weights realizes $C\log p$ as orbit
periods for all primes. Any exact construction must leave the finite-metric-graph,
additive-travel-time regime. [evidence type A]
\end{theorem}

The escape routes are an infinite graph, a non-additive or state-dependent weight, or an
operator whose spectrum encodes arithmetic without the classical lengths literally being
$\log p$. We note that Hilbert--P\'olya never required anyone to build racetracks of
length $\log p$; the earlier possibility of an explicit optical construction of the
primes is closed, not the operator route.

\paragraph{What survives: counting, not equality.} A weaker target does survive, and a
toy computation shows both why it looks impressive and why it is not. A bouquet of three
loops $a,b,c$ with lengths $\ell_a,\ell_b,\ell_c$ generates primitive orbits as
primitive cyclic words, with periods forced to be
$L(w) = n_a\ell_a + n_b\ell_b + n_c\ell_c$; no prime is supplied to the model. Fitting
the three lengths to $\log p_1,\dots,\log p_N$ (enumerating primitive necklaces to
word-length 11, 25{,}486 orbits) reproduces:
\begin{center}
\begin{tabular}{rrrrr}
\toprule
$N$ & fitted $(\ell_a,\ell_b,\ell_c)$ & RMS rel.\ error & worst & entropy $h$ \\
\midrule
20 & $(0.7229,\,1.1274,\,1.6602)$ & $2.13\%$ & $4.92\%$ & $0.99937$ \\
30 & $(0.7190,\,1.1279,\,1.6672)$ & $1.80\%$ & $5.09\%$ & $0.99972$ \\
50 & $(0.7185,\,1.1309,\,1.6672)$ & $1.49\%$ & $4.96\%$ & $0.99903$ \\
\bottomrule
\end{tabular}
\end{center}
For a suspension on a three-loop bouquet the topological entropy solves
$e^{-h\ell_a}+e^{-h\ell_b}+e^{-h\ell_c}=1$, and the fitted lengths land at
$h \approx 0.999$. This has been read as an independent discovery that the system
self-organises to the entropy required by the prime number theorem under logarithmic
coordinates. It is not.

\begin{proposition}[$h=1$ is forced by the objective]\label{prop:hforced}
Suppose a hyperbolic flow satisfies $\pi_{\mathrm{orb}}(T) \sim e^{hT}/(hT)$ and its
primitive periods satisfy $L_n \sim \log p_n$. Then $h = 1$. [evidence type A]
\end{proposition}
\begin{proof}
Evaluate the counting function at $T=\log p_n$. Its left side is $n$; its right side is
$e^{h\log p_n}/(h\log p_n) = p_n^{h}/(h\log p_n)$. Since $n \sim p_n/\log p_n$, agreement
requires $p_n^{h}/h \sim p_n$, i.e.\ $p_n^{h-1}\sim h$, which forces $h=1$.
\end{proof}

So the optimiser was handed the answer: matching $L_n$ to $\log p_n$ forces the counting
law to agree asymptotically, and the counting law is what pins $h$. The $0.1\%$
departure from unity is the finite window and the fitting residual, not signal. The
methodological point is the useful one: leading-order agreement is automatic and
therefore worthless as evidence, so any advance in this direction has to live in the
fluctuation term.

\paragraph{Where the zeros would have to be.} That fluctuation term is the only version
of this programme with a chance of being interesting. Normalising $h=1$, one would have
to match
\begin{equation}
\Delta_{\mathrm{opt}}(T) = \pi_{\mathrm{orb}}(T) - \operatorname{Li}(e^T)
\quad\text{to}\quad
\Delta_{\mathrm{prime}}(T) = \pi(e^T) - \operatorname{Li}(e^T),
\end{equation}
not the growth rate but the error. Each orbit would additionally need its linearized
Poincar\'e return map $P_\gamma$ on the four-dimensional contact plane $\ker\alpha$, and
a Maslov-type index---the data a Gutzwiller-type trace formula consumes. None of that is
constructed here.

\paragraph{Net position.} What this appendix establishes is narrow and firm. The
closed-ray obstruction of \S\ref{sec:extensions} is a periodic-orbit statement wearing a
memory-layout costume; the dynamical zeta function is a natural and harmless packaging of
closed-ray delays; the finite-graph route to arithmetic is provably closed; and
leading-order agreement with the prime-log counting law is automatic. The programme is
not refuted, but its easy version is, and the remaining version is a fluctuation problem
whose entire content is the error term.

\section*{Data and code availability}
The released artifacts accompany this manuscript: the synthetic ray generator and
deterministic split rule (\texttt{code/generate.py}), the uniform and chart-legal DP
partition code with 30-seed evaluation (\texttt{code/partition.py}), the per-seed
results (\texttt{code/results\_exp5\_seeds.csv}), the budget-scaling, control, analytic-partition
and chart-legality analysis behind Tables in \S\ref{sec:experiments}
(\texttt{code/scaling.py}, \texttt{code/results\_exp5\_scaling.csv}), the Ada/SPARK encoder with its
exhaustive runtime test (\texttt{code/chart\_encoder.*}, \texttt{code/test\_encoder.adb}),
and the GNATprove transcript (\texttt{code/gnatprove-transcript.txt}). The two
calibration constants of the reference generator are documented in its source; the
chart table is the immutable table of Experiment~4. An independent audit of these
artifacts, including the corrections it prompted, is released as
\texttt{code/AUDIT.md}.

\bibliographystyle{unsrtnat}
\bibliography{refs}

\end{document}
